\CWHeader{Лабораторная работа \textnumero 5}

\CWProblem{
Найти в заранее известном тексте поступающие на вход образцы с использованием суффиксного массива.
}

\textbf{Формат ввода}
Текст располагается на первой строке, затем, до конца файла, следуют строки с образцами.

\textbf{Формат вывода}
Для каждого образца, найденного в тексте, нужно распечатать строчку, начинающуюся с последовательного номера этого образца и двоеточия, за которым, через запятую, нужно перечислить номера позиций, где встречается образец в порядке возрастания.

\textbf{Примечание}
Можно делать как через суффиксное дерево, так и через сортировку.

\pagebreak